\documentclass[12pt letter]{report}
\input{./template/preamble}
\input{./template/macros}
\input{./template/letterfonts}

\title{\Huge{Wireless and Mobile Networks}}
\author{\huge{Madiba Hudson-Quansah}}
\date{}
\usepackage{enumitem}
\usepackage{parskip}

\setcounter{tocdepth}{4}
\setcounter{secnumdepth}{4}

\begin{document}
\maketitle
\newpage
\pdfbookmark[section]{\contentsname}{too}
\tableofcontents
\pagebreak

\chapter{Introduction}

The following are elements of a wireless network:
\begin{description}
  \item[Wireless Hosts]  - End-system devices that run applications,
    for example smartphones, tablets, laptops, and IoT devices.
  \item[Wireless links] - A host connects to a base station (BS) or
    to another wireless host, via a \textbf{wireless communication
    link}. Different wireless link technologies have different
    transmission rates and can transmit over different distances.
  \item[Base Station] - A base station is responsible for sending and
    receiving data to and from a wireless host that is associated
    with that base station. A base station is often responsible for
    coordination the transmission of multiple wireless hosts with
    which it is associated with. Association in this context means that:
    \begin{enumerate}
      \item The host is within the wireless communication distance of
        the base station
      \item The host uses that base station to relay data between it,
        the host, and the larger network.
    \end{enumerate}
    Cell towers and access points are examples of base stations.
    Hosts associated with a base station are often referred to as
    operating in \textbf{infrastructure mode}, since all traditional
    network services are provided by the network to which a host is
    connected via the base station. In \textbf{ad hoc networks},
    wireless hosts have no such infrastructure with which to connect.

    When a mobile host moves beyond the range of one base station and
    into the range of another, it will change its point of attachment
    into the larger network, this is referred to as \textbf{handoff / handover}.
  \item[Network Infrastructure] - The larger network with which a
    wireless host may wish to communicate.
\end{description}

At the highest level we can classify wireless networks based on two criteria:
\begin{itemize}
  \item Whether a packet in the wireless network crosses exactly one
    wireless hop or multiple wireless hops.
  \item Whether there is infrastructure such as a base station in the network.
\end{itemize}

\chapter{Wireless Links and Network Characteristics}

Wireless links differ from wired links in several ways:
\begin{description}
  \item[Decreasing Signal Strength]  - Even in free space the signal
    will disperse resulting in decreased signal strength, also called
    \textbf{path loss}, as the distance between sender and receiver increases.
  \item[Interference from other sources] - Radio sources transmitting
    in the same frequency band will interfere with each other.
  \item[Multipath Propagation] - Occurs when portions of the
    electromagnetic wave reflect off objects and the ground, taking
    paths of different lengths between a sender and receiver. This
    results in the blurring of the received signal at the receiver.
\end{description}

The \textbf{signal-to-noise ratio (SNR)} is a relative measure of the
strength of the received signal and this noise. The SNR is measured
in units of decibels (dB). A larger SNR makes it easier for the
receiver to extract the transmitted signal from the background noise.

The \textbf{bit error rate (BER)} is the probability that a
transmitted bit is received in error at the receiver.

There are several physical-layer characteristics that are important
to understanding higher-layer wireless communication protocols:
\begin{description}[style=multiline, leftmargin=7cm]
  \item[For a given modulation scheme, the higher the SNR, the lower
    the BER]  - Since a sender can increase SNR by increasing its
    transmission power, a sender can decrease the probability that a
    frame is received in error by increasing its transmission power.
  \item[For a given SNR, a modulation technique with a higher bit
    transmission rate will have a higher BER] - This is because a
    modulation technique with a higher bit transmission rate will
    transmit more bits per symbol, and thus the probability that at
    least one of those bits is received in error will be higher.
  \item[Dynamic Selection of the physical-layer modulation technique
      can be used to adapt the modulation technique to channel
    conditions] - The SNR, and hence BER, may change as a result of
    mobility or due to changes in the environment. By dynamically
    selecting the modulation technique, a sender can adapt to changing
    channel conditions, for example by switching to a modulation technique
    with a lower bit transmission rate when the SNR is low, and switching
    to a modulation technique with a higher bit transmission rate when the
    SNR is high.
\end{description}

The \textbf{Hidden Terminal Problem}, occurs when physical
obstructions in the environment prevent two wireless hosts from
sensing each other's transmissions, even though they are both within
the communication range of a common base station. This can lead to
collisions at the base station when both hosts transmit
simultaneously, as they are unaware of each other's transmissions.

\textbf{Fading} is a phenomenon that occurs when a signal's strength
as it propagates through the wireless medium experiences variations
due to factors such as multipath propagation, shadowing, and
atmospheric conditions. Fading can lead to fluctuations in the
received signal strength, which can affect the reliability of
wireless communication.

\section{Code Division Multiple Access (CDMA)}

\dfn{Chipping Rate}{
  The rate at which the code changes in a CDMA protocol, which is
  much faster than the bit transmission rate.
}

In a Code Division Multiple Access (CDMA) protocol each bit being
sent is encoded by multiplying the bit by a signal (the code) that
changes at a much faster rate, i.e. \textbf{chipping rate}, than the
original sequence of data bits.

Supposing that the rate at which original data bits reach the CDMA
defines the unit of time, i.e. each original data bit to be
transmitted requires a one-bit slot time. Let $d_i$ be the value of
the data bit for the $i$-th bit slot. We represent a data bit with a
0 value as -1. Each bit slot is further subdivided into $M$ mini-slots.

CDMA works by assigning a unique code to each sender, and encoding
the data bits by multiplying them with the sender's code. The
receiver can then decode the received signal by multiplying it with
the code of the sender it wishes to receive data from. This allows
multiple senders to transmit simultaneously over the same frequency
band without interfering with each other, as long as their codes are
orthogonal to each other.

For example suppose that there are two senders, A and B, with codes
$c_A$ and $c_B$ respectively. If sender A wants to transmit a data
bit $d_A$ and sender B wants to transmit a data bit $d_B$ at the same
time, they will both go through the following process:
\begin{enumerate}
  \item Sender A encodes its data bit by multiplying it with its
    code: $s_A = d_A \cdot c_A$.
  \item Sender B encodes its data bit by multiplying it with its
    code: $s_B = d_B \cdot c_B$.
  \item Both senders transmit their encoded signals simultaneously
    over the same frequency band, resulting in a combined signal at
    the receiver: $s = s_A + s_B$.
  \item The receiver can decode the signal from sender A by
    multiplying the combined signal with sender A's code: $s \cdot
    c_A = (s_A + s_B) \ cdot c_A = s_A \cdot c_A + s_B \cdot c_A$.
    Since $c_A$ and $c_B$ are orthogonal, the term $s_B \cdot c_A$
    will be zero, allowing the receiver to recover sender A's data
    bit: $s \cdot c_A = s_A \cdot c_A = d_A \cdot (c_A \cdot c_A) = d_A$.
  \item Similarly, the receiver can decode the signal from sender B by
    multiplying the combined signal with sender B's code: $s \cdot
    c_B = (s_A + s_B) \cdot c_B = s_A \cdot c_B + s_B \cdot c_B$.
    Again, since $c_A$ and $c_B$ are orthogonal, the term $s_A \cdot
    c_B$ will be zero, allowing the receiver to recover sender B's data
    bit: $s \cdot c_B = s_B \cdot c_B = d_B \cdot (c_B \cdot c_B) = d_B$.
\end{enumerate}

\chapter{Wireless Fidelity (WiFi): 802.11 Wireless LANs}

\section{The 802.11 Wireless LAN Architecture}

\subsection{Channels and Association}

\section{The 802.11 MAC Protocol}

\subsection{Dealing with Hidden Terminals: Request to Send (RTS) and
Clear to Send (CTS)}

\subsection{Using 802.11 as a Point-to-Point Link}

\section{The IEEE 802.11 Frame}

\subsection{Payload and CRC Fields}

\subsection{Address Fields}

\subsection{Sequence Number, Duration and Frame Control Fields}

\section{Mobility in the same IP Subnet}

\section{Advanced Features in 802.11}

\subsection{802.11 Rate Adaptation}

\subsection{Power Management}

\section{Personal Area Networks: Bluetooth}

\chapter{Cellular Networks: 4th Generation and 5th Generation}

Cellular refers to the fact that the region covered by a cellular
network is divided into a number of geographic coverage areas called
\textbf{cells}. Each cell contains a \textbf{base station} that
transmits and receives signals to and from \textbf{mobile devices}
currently in its cell.

\section{4th Generation (4G) Long Term Evolution (LTE) Cellular Networks}

Fourth Generation (4G) networks implement the Long-Term Evolution
(LTE) standard. The network broadly divides into the radio network at
the cellular network's edge and the core network. All network
elements communicate with each other using IP. The elements that make
up a 4G LTE network include:
\begin{description}
  \item[Mobile Device]  - A smartphone, tablet, laptop, or IoT device
    that connects into a cellular carrier's network. The mobile
    device has a globally unique 64-bit identifier called the
    \textbf{International Mobile Subscriber Identity (IMSI)} which is
    stored on its \textbf{Subscriber Identity Module (SIM) } card.
    The IMSI identifies the subscriber in the worldwide cellular
    carrier network system, including the country and home cellular
    carrier network to which the subscriber belongs. The mobile
    device is referred to as \textbf{User Equipment (UE)} in the LTE standard.
  \item[Base Station] - Base stations sit at the edge of the
    carrier's network and is responsible for managing the wireless
    radio resources and the mobile devices within its coverage area.
    Mobile devices will interact with a base station to attach to the
    carrier's network. Cellular base station functions are comparable
    to those of APs in wireless LANs, bit also serve other important
    roles such as creating device-specific IP tunnels from the mobile
    device to gateways and interacting amongst themselves to handle
    device mobility among cells.
  \item[Home Subscriber Server (HSS)] - A control-plane database that
    stores information about the mobile devices for which the HSS's
    network is their home network. Can be used with MME for device
    authentication
  \item[Serving Gateway (S-GW), Packet Data Network Gateway (P-GW),
    and other network routers] - The S-GW and P-GW are two routers
    that lie on the data path between the mobile device and the
    Internet. The P-GW also provides NAT IP addresses to mobile
    devices and performs NAT functions. The P-GW is the last LTE
    element that a datagram originating at the mobile device will
    pass before entering the larger Internet. A cellular carrier's
    all-IP core will have additional routers whose role is similar to
    that of traditional IP routers, i.e. to forward IP datagrams
    among themselves along paths that will typically terminate at
    elements of the LTE core network.
  \item[Mobility Management Entity (MME)] - A control-plane element
    that plays an important role in authenticating a device wanting
    to connect into its network. It also sets up the tunnels on the
    data path from/to the device and the P-GW, and maintains
    information about an active mobile device's cell location within
    the carrier's cellular network.
\end{description}

\section{LTE Protocol Stacks}

\section{LTE Radio Access Network}

\section{Additional LTE Functions: Network Attachment and Power Management}

\subsection{Network Attachment}

\subsection{Power Management: Sleep Modes}

\section{The Global Cellular Network: A Network of Networks}

\section{5th Generation (5G) Cellular Networks}

\subsection{5G and Millimetre Wave Frequencies}

\subsection{5G Core Network}

\chapter{Mobility Management: Principles}

\section{Device Mobility: A Network Layer Perspective}

\section{Home Networks and Roaming on Visited Networks}

\section{Direct and Indirect Routing to/from a Mobile Device}

\subsection{Leveraging the Existing IP Address Infrastructure}

\subsection{Indirect Routing to a Mobile Device}

\subsection{Direct Routing to a Mobile Device}

\chapter{Mobility Management in Practice}

\section{Mobility Management in 4G/5G Networks}

\section{Mobile IP}

\chapter{Wireless and Mobility: Impact on Higher-Layer Protocols}

\end{document}
